\chapter{Resultados y evaluación}

Este capítulo define el protocolo con el que se realizará la evaluación definitiva del prototipo. Las ejecuciones realizadas durante el desarrollo se utilizaron para depurar el sistema, pero no se presentan como resultados finales: la comparación definitiva debe realizarse con una batería fijada previamente, los mismos datos locales y criterios de revisión homogéneos.

\section{Objetivo de la evaluación}

El objetivo no es demostrar que el sistema prediga el mercado, sino comprobar si la arquitectura interpreta consultas financieras históricas, genera código Python adecuado, valida su seguridad, ejecuta métricas trazables y redacta conclusiones comprensibles.

El proyecto contempla tres perfiles de ejecución:

\begin{itemize}
    \item Groq con \nolinkurl{llama-3.3-70b-versatile} \cite{groqllama33docs}.
    \item Gemini con \nolinkurl{gemini-2.5-flash} \cite{google2025gemini25flashdocs}.
    \item Perfil técnico \texttt{university}, correspondiente a \nolinkurl{openai/gpt-oss-20b} \cite{openai2025gptossmodelcard} servido mediante vLLM en el entorno universitario.
\end{itemize}

En todos los casos el workflow usa el modelo configurado para interpretación de consulta, generación de código e interpretación final. Si falta una clave, el proveedor no está disponible, el código se rechaza o la ejecución falla, el caso debe quedar registrado como error controlado.

\section{Métricas y criterios}

La evaluación técnica registrará:

\begin{itemize}
    \item \textbf{Casos completados}: consultas válidas que terminan con estado \texttt{completed}.
    \item \textbf{Errores controlados}: consultas inválidas o fallidas que terminan con explicación.
    \item \textbf{Errores LLM}: problemas de configuración, red, cuota, disponibilidad, contexto o formato.
    \item \textbf{Código rechazado}: scripts que no superan el contrato de seguridad.
    \item \textbf{Errores de ejecución}: scripts válidos que fallan al calcular o serializar la salida.
    \item \textbf{Violaciones de recomendación de inversión}: presencia de recomendaciones o predicciones fuera del alcance del MVP.
    \item \textbf{Tiempo de ejecución}: duración observada por caso y proveedor.
\end{itemize}

Estas métricas deben complementarse con revisión humana. La completitud técnica no garantiza por sí sola que el análisis sea útil. La revisión cualitativa valora comprensión de la consulta, selección de métricas, suficiencia del código generado, adecuación de tablas o gráficas, exactitud respecto a los datos ejecutados, claridad del texto final, reconocimiento de limitaciones y ausencia de recomendaciones de inversión.

Esta decisión es coherente con la necesidad de explicar cómo se calcula y presenta el rendimiento histórico \cite{investorgov2022performanceclaims}, y con la separación entre retorno, riesgo, volatilidad y correlación como conceptos de análisis histórico \cite{investorgov2026risk,cfainstitute2026portfolioriskreturn}.

La rúbrica se define en \nolinkurl{docs/evaluacion_cualitativa_llm.md}. En ella se separan criterios del primer agente, como selección de métricas y generación de datos suficientes, y criterios del segundo agente, como claridad, prudencia y fidelidad a la salida ejecutada.

\section{Niveles cualitativos de salida}

La evaluación se estructura en tres niveles:

\begin{itemize}
    \item \textbf{Nivel A}: consultas simples sin requisitos de presentación. La respuesta esperada es breve y debe incluir métricas básicas suficientes.
    \item \textbf{Nivel B}: consultas que piden un análisis más completo. La respuesta esperada incluye texto, tabla de métricas y una gráfica principal o datos estructurados para generarla.
    \item \textbf{Nivel C}: consultas profesionales o detalladas. La respuesta esperada incluye varias salidas estructuradas, por ejemplo tabla, evolución normalizada, drawdown, clasificación mensual, desglose de riesgo o conclusión ejecutiva.
\end{itemize}

Para series temporales se priorizan gráficos de línea y, cuando se comparan activos con escalas distintas, una evolución normalizada permite comparar la tendencia sin confundir niveles de precio absolutos. También se evita abusar de demasiadas series o ejes dobles, porque pueden dificultar la lectura \cite{atlassianlinecharts}. En consultas de riesgo se considera especialmente útil el drawdown, ya que resume caídas históricas desde máximos previos y se entiende con facilidad \cite{cfainstitute2018drawdown}.

\begin{table}[htbp]
\centering
\small
\begin{tabularx}{\textwidth}{l X X}
\toprule
\textbf{Nivel} & \textbf{Uso previsto} & \textbf{Salida esperada} \\
\midrule
A & Query simple sin formato específico & Texto breve y métricas básicas en tabla o bloque estructurado. \\
B & Análisis más completo o comparativa clara & Texto, tabla de métricas y una gráfica principal o datos para representarla. \\
C & Petición profesional, detallada o con varias salidas & Texto, tablas, varias visualizaciones o datos estructurados, conclusión y limitaciones. \\
\bottomrule
\end{tabularx}
\caption{Niveles cualitativos de profundidad del análisis.}
\label{tab:niveles_cualitativos_analisis}
\end{table}

\section{Batería progresiva}

La batería cualitativa incluye consultas A/B/C sobre crecimiento, comparación de activos, visión descriptiva, retornos, riesgo histórico y análisis técnico. También se mantiene una query de estrés visual que obliga al sistema a respetar cuatro bloques de salida:

\begin{quote}
Compara QQQ y SPY durante 2024 como si fuera un informe para un cliente no técnico. Quiero que me muestres los datos exactamente en cuatro bloques: 1) una tabla resumen con rentabilidad total, volatilidad anualizada, máximo drawdown, mejor mes y peor mes; 2) una tabla de clasificación mensual indicando qué activo ganó cada mes; 3) datos para una gráfica normalizada base 100 de ambos activos; 4) datos para una gráfica de drawdown. Cierra con una conclusión clara, sin recomendar comprar ni vender.
\end{quote}

Esta consulta comprueba simultáneamente comprensión, selección de métricas, generación de código, control de formato, datos para visualización e interpretación final. La batería completa y las familias técnicas complementarias se describen en \nolinkurl{docs/planificacion_evaluacion_tfm.md}.

\section{Protocolo de ejecución definitiva}

La evaluación final seguirá los siguientes pasos:

\begin{enumerate}
    \item Fijar la batería exacta de consultas y los CSV locales utilizados.
    \item Ejecutar la misma batería con cada perfil LLM seleccionado.
    \item Guardar por caso el plan, el código, la salida estructurada, los avisos, el estado y el tiempo.
    \item Revisar cada respuesta con la plantilla manual definida previamente.
    \item Separar fallos de proveedor, fallos de codegen, rechazos de seguridad, fallos de ejecución y defectos de interpretación.
    \item Presentar la comparación como evidencia observada bajo condiciones concretas, no como clasificación estadística absoluta.
\end{enumerate}

\section{Resultados pendientes}

Los resultados comparativos definitivos se incorporarán después de ejecutar la batería cerrada. Hasta entonces no se presentan porcentajes de completitud, latencias por proveedor ni ejemplos numéricos como conclusiones finales del TFM.

Los informes Markdown y JSON producidos por scripts de evaluación se consideran artefactos temporales mientras no se haya seleccionado la pasada definitiva. Esta separación evita confundir registros de depuración con evidencia final.

\section{Limitaciones metodológicas}

La calidad de las respuestas LLM requiere revisión humana y las cifras de latencia dependen de red, cuota y carga del proveedor. Además, una batería reducida no permite afirmar robustez estadística. La comparación debe entenderse como exploratoria, trazable y revisable.

El sistema trabaja con datos históricos locales. Su objetivo es describir métricas observadas, no producir asesoramiento financiero, incorporar noticias externas ni predecir el comportamiento futuro de los activos.

\section{Síntesis}

La evaluación propuesta combina criterios técnicos y cualitativos. Su aportación principal consiste en revisar por separado interpretación, generación de código, validación, ejecución e interpretación final. Esta estructura permite analizar no solo si el workflow termina, sino también si responde con datos suficientes, comprensibles y trazables.
